\documentclass[11pt]{article}

\usepackage[utf8]{inputenc}
\usepackage[T1]{fontenc}
\usepackage[margin=1in]{geometry}
\usepackage{amsmath}
\usepackage{amsfonts}
\usepackage{amssymb}
\usepackage{booktabs}
\usepackage{hyperref}
\usepackage{xcolor}
\usepackage{microtype}
\usepackage{parskip}
\usepackage{enumitem}

\hypersetup{
  colorlinks=true,
  linkcolor=blue!50!black,
  urlcolor=blue!50!black,
  citecolor=blue!50!black
}

\title{Bottleneck of GNN connectome inference under \\ measurement noise --- synthesis and implementation plan}
\author{}
\date{}

\begin{document}

\maketitle

\tableofcontents
\bigskip\hrule\bigskip

\section{Where the wall is, mathematically}

The flyvis ODE is \emph{nonlinear in voltages} but \emph{linear in circuit
parameters} $\theta = (\tau_i, W_{ij}, V^{\text{rest}}_i)$. Stacking
the simulator update rule across $(N\,T)$ samples gives an over-determined
linear system $A\theta = b$ where the design matrix $A$ depends only on the
observed voltages and stimuli (\texttt{neurips.tex:217}, app.\ \S B). The
inverse problem is therefore well-posed up to the null space of $A$, and the
identifiability of $\theta$ is purely a question of whether $A^\top A$ is
invertible.

It is not. Two structural facts kill recovery before a single noise photon
enters:

\subsection{Continuous degeneracy in $W$ (geometric symmetry)}

In the central FlyVis crop ($N{=}13{,}741$ neurons, $E{=}434{,}122$ edges) the
stimulus-driven Hessian $H = A^\top A$ has $\dim\ker(H) = 308{,}160$ ---
roughly $71\%$ of the edges. The reason is structural: same-type neurons
across different retinotopic columns receive \emph{stimulus-shifted copies}
of the same input. If $W_{ij}$ is the weight feeding neuron $i$ from neuron
$j$, and $W_{i'j'}$ is the analogous weight one column over, then the data
$\{v(t)\}$ supports a continuous one-parameter family of perturbations
$W_{ij} \mapsto W_{ij} + \delta$, $W_{i'j'} \mapsto W_{i'j'} - \delta$ that
leave the predicted dynamics unchanged --- \emph{sum-zero} edge perturbations
within columnar equivalence classes. Empirically the appendix shows $W$ can
be perturbed by $\geq 8\times$ the group mean while rollout Pearson $r > 0.94$
(\texttt{neurips.tex:1061}, Table~$\lambda_8$).

\subsection{Scale degeneracy $W \leftrightarrow f_\theta$ (architectural symmetry)}

The GNN forward model is
\[
\frac{d\hat v_i}{dt} = f_\theta\!\Big(v_i,\, \mathbf{a}_i,\,
\underbrace{\textstyle\sum_j \widehat W_{ij}\, g_\phi(v_j,\mathbf{a}_j)}_{=\,m_i},\,
I_i\Big).
\]
For \emph{any} positive constant $k$, the substitution
$(\widehat W,\, f_\theta(\cdot,\cdot,m,\cdot)) \mapsto (k\widehat W,\, f_\theta(\cdot,\cdot,m/k,\cdot))$
leaves $d\hat v/dt$ pointwise unchanged, because $k\widehat W \cdot g_\phi$
multiplies $m$ by $k$ and the message branch of $f_\theta$ divides it back by $k$
on the way out. This is an \emph{exact} continuous symmetry of the
architecture. The MLP $f_\theta$ has more than enough capacity to absorb a
factor of $\sim 2$ on its message branch, and the data alone cannot tell which
$k$ is the ``right'' one.

The consequence: the recovered $\widehat W$ is determined by the data
\emph{only up to an unknown global scale}. The regression slope of
$\widehat W$ vs.\ $W_{\text{true}}$ depends on which point $k^\star$ the
optimizer landed on. In our baseline that point is $k^\star \approx 0.51$
(\texttt{instruction:120--138}).

\subsection{What it means for measurement noise to be ``load-bearing''}

\textbf{Definition.} A symmetry is ``load-bearing'' if the optimizer's choice
within the symmetry orbit is \emph{determined} by the noise rather than by the
signal. In the noise-free case, all points on the orbit yield the same loss,
and the optimizer's landing point is set by initialisation, regularization,
and gradient-flow geometry --- not by the data. A modest L1 penalty on
$\widehat W$ is enough to nudge the system to the natural-scale point
$k^\star \approx 1$, recovering $W$ R$^2 \approx 0.97$.

Add measurement noise. The training signal is now a \emph{noisy} finite
difference $\widetilde y_i(t) = (v_i(t{+}1) - v_i(t))/\Delta t + \eta_i(t)$
where $\eta_i$ has standard deviation $\gamma\sqrt{2}/\Delta t$. At
$\gamma=0.10$ and $\Delta t = 0.02$ this is $\eta\text{-std} \approx 7.07$ ---
the same order of magnitude as the clean derivative itself
(\texttt{instruction:34}). The MSE loss now has two contributions:
a signal term that is invariant along the orbit, and a noise term
$\propto \mathbb E\|\eta\|^2 \sim k^2 \cdot \mathrm{Var}(\eta_{\text{eff}})$
that is \emph{not}, because the message-branch Jacobian
$\partial f_\theta / \partial m \propto 1/k$ scales the variance the noise sees
on its way through the network. The optimizer minimises this term by
\emph{shrinking} $\widehat W$ (smaller $k$) and inflating $f_\theta$'s
message branch correspondingly.

\textbf{The slope-shrunk $W$ is therefore the path of least resistance
under noise:} of all points on the symmetry orbit, the shrunken one is the
local minimum of the noise penalty. The network has no incentive to leave it
because every voltage-level loss is constant along the orbit. This is what
``load-bearing'' means here: the symmetry is silent in the clean case but
\emph{actively selects the wrong $W$} as soon as noise enters the gradient.

\subsection{What ``no information orthogonal to the symmetry'' means}

A loss function $\mathcal L$ is \emph{symmetry-blind} if
$\mathcal L(\widehat W, f_\theta) = \mathcal L(k\widehat W, f_\theta(\cdot,m/k,\cdot))$
for all $k>0$. Any voltage-level loss --- one-step MSE, multi-step rollout
MSE, smoothness on $v$, EMA-denoised $v$, derivative-SNR-weighted MSE,
voltage-coherence terms across neurons --- is symmetry-blind by construction,
because all of these are functions of $v$ and $\hat v$, and the symmetry
\emph{by definition} preserves $\hat v$ pointwise.

Adding more such losses cannot break the degeneracy. Geometrically: the
symmetry traces out a one-dimensional orbit in parameter space; any
symmetry-blind loss is constant along that orbit; the gradient of such a
loss has zero component in the orbit direction. Information ``orthogonal to
the symmetry'' means a loss term whose gradient \emph{has} a non-zero
component in that direction --- a constraint that distinguishes $k$ from
$1/k$. Examples: $g_\phi(v_\star,\mathbf a) = v_\star$ (anchors $g_\phi$ at a
reference voltage), $\partial f_\theta/\partial m = 1$ (anchors the
message-branch Jacobian), $\sum_{ij}|W_{ij}| \cdot \mathrm{sign}(W^{\text{Dale}}_{ij})$
(Dale's law: signs of edges by presynaptic type), or a typed equivariance
constraint that ties $W_{ij}$ values across columns of the same
presynaptic-postsynaptic type pair.

This is why the LLM exploration's voltage-side blocks (denoising, coherence,
EMA --- all symmetry-blind) reverted: they cannot, even in principle, move
the optimizer off the slope-shrunk minimum.

\section{Empirical evidence from the three explorations}

\begin{table}[h]
\centering
\small
\begin{tabular}{lll}
\toprule
\textbf{Regime} & \textbf{Best $W$ R$^2$} & \textbf{Notes} \\
\midrule
\texttt{noise\_005} ($\sigma=0.05$ only)
  & 0.983 (CV 0.21\%)
  & Recurrent + msg\_diff=100 + g\_phi\_L1=0; near-oracle \\
\texttt{noise\_005\_010} ($\sigma=0.05,\, \gamma=0.10$)
  & 0.7457 $\rightarrow$ 0.802
  & Stuck below 0.78 oracle; $W$ slope $\approx 0.51$ \\
\texttt{noise\_005\_010\_rc} (recurrent fine-tune)
  & 0.802 (5-seed)
  & Bimodal convergence --- Batch~32 collapsed (0.07--0.17) \\
\texttt{removed\_pc\_20} (structural perturbation, no $\gamma$)
  & 0.887
  & Edge removal is \emph{easier} than measurement noise \\
\bottomrule
\end{tabular}
\end{table}

Adding $\gamma=0.10$ costs $\sim 0.18$ absolute $W$ R$^2$. Five LLM-explored
anchor families that touch only voltage-level signals (denoising, EMA,
coherence, $g_\phi$ 10-pt identity, $g_\phi$ zero-intercept) \textbf{all
reverted}: as just argued, they re-process the same noise-corrupted observable
with symmetry-blind losses, adding no information that distinguishes the
correct $k$ from any other (\texttt{instruction:120--209}).

\section{Why the GNN inherits this from the loss, not the architecture}

The Known-ODE oracle (perfect ReLU, only $W$, $\tau$, $V^{\text{rest}}$ free,
\emph{no} message-branch MLP and therefore \emph{no} scale degeneracy)
plateaus at $W$ R$^2 \approx 0.78$--$0.80$ at $\gamma=0.10$
(\texttt{neurips.tex:1393}). The GNN at 0.745 sits in the same basin: the
gap to oracle is architectural (the MLP gives it the scale degeneracy that
the oracle lacks), but the gap from oracle to noise-free
(0.78 $\rightarrow$ 0.965) is \emph{genuinely loss-limited}. Even with no
parametric symmetry, the pointwise MSE on $(v[t{+}1]-v[t])/\Delta t$ is
half-drowned by derivative noise of std~$\sim 7$, so the conditioning of
$A^\top A$ collapses regardless of the architectural choice. The neurips
paper concedes the point and offers no fix (line 584).

\medskip
\noindent\textbf{Two distinct attacks are needed:}
\begin{enumerate}[label=(\Alph*),leftmargin=*]
  \item \textbf{Reformulate the loss} so it is no longer pointwise on a
    finite difference (attacks the noise wall --- closes the
    $0.78 \rightarrow 0.965$ gap that even the oracle cannot close).
  \item \textbf{Add a symmetry-breaking anchor} so the optimizer cannot rest
    at the slope-shrunk point (attacks the architectural gap ---
    closes $0.745 \rightarrow 0.78$).
\end{enumerate}
The rest of this document is about (A), with implementation guided by SPEND.

\section{What SPEND (Ding et al.\ 2025) actually offers}

SPEND solves a \emph{different} problem: spatially-correlated,
spectrally-varied noise in hyperspectral SRS/MIP imaging, denoised via
Noise2Noise (N2N) self-permutation along an axis where the noise is least
correlated. We inspected the reference implementation
(\texttt{/workspace/connectome-gnn/papers/SPEND/}); the relevant facts are:

\begin{itemize}[leftmargin=*]
  \item \textbf{Permutation pair construction
    (\texttt{datasplit\_with\_aug\_choose.py:75--97}).}
    Literal stride-2 split + concatenate along the chosen axis $\omega$:
\begin{verbatim}
    even = img[::2, :, :]
    odd  = img[1::2, :, :]
    input_stack  = concat(even, odd)
    target_stack = concat(odd,  even)
\end{verbatim}
    Both stacks are noisy. Both have \emph{the same underlying signal} (up
    to a one-frame shift along $\omega$) but \emph{independent noise
    realisations}. That is the definition of an N2N pair.
  \item \textbf{Loss (\texttt{losses.py:37--48}).}
    Pure MAE between network output and target; no masking, no perceptual
    term.
  \item \textbf{Model.} 3D U-Net via CSBDeep, depth~4, base channels 32, no
    batch norm, residual skip on the output, linear head.
  \item \textbf{Training.} Adam, LR $4\times10^{-5}$, 100 epochs, batch~16.
    Augmentation is flips only (no rotation), applied at preprocessing time.
  \item \textbf{Inference (\texttt{3\_prediction\_SPEND.py}).} The trained
    network is applied to the \emph{original} (un-permuted) sequence; the
    permutation is a training-time trick only.
  \item \textbf{No explicit noise model.} The network is purely supervised
    on N2N pairs --- the only ``model'' of the noise is the implicit
    assumption that it decorrelates along the chosen permutation axis.
\end{itemize}

\section{Why SPEND-style training applies even though our noise is i.i.d.}

This is the key sceptic-pass. Our measurement noise is added independently
per $(\text{neuron}, \text{frame})$:
$v^{\text{obs}}_i(t) = v^{\text{clean}}_i(t) + \gamma\, \xi_i(t)$,
$\xi_i(t) \overset{\text{iid}}{\sim} \mathcal N(0,1)$.
SPEND was \emph{designed} for non-independent noise. Why use a
non-independent-noise method on independent noise?

\subsection{The Noise2Noise theorem does not require non-independent noise}

The original Noise2Noise result (Lehtinen et~al.\ 2018) shows that, given
any two noisy observations $\widetilde x_1 = s + \eta_1$ and
$\widetilde x_2 = s + \eta_2$ with $\eta_1, \eta_2$ \emph{independent}
zero-mean noise, the network trained with
$\mathcal L = \mathbb E \|f(\widetilde x_1) - \widetilde x_2\|^2$
converges to $\mathbb E[\widetilde x_2 \mid \widetilde x_1] =
\mathbb E[s \mid \widetilde x_1] + 0 = $ posterior mean of $s$. The proof
hinges on $\mathbb E[\eta_2 \mid \widetilde x_1] = 0$, which holds whenever
$\eta_1 \perp \eta_2$. \emph{i.i.d.\ Gaussian noise satisfies this trivially}
--- in fact it was the canonical case in the original paper.

SPEND's contribution was to show that even when noise is \emph{not} i.i.d.\
across all axes (it is correlated along the fast-scan axis in SRS), there is
\emph{still} an axis along which it decorrelates (the slow-scan or spectral
axis), and permuting along that axis gives valid N2N pairs. SPEND extends
N2N to a harder regime; it does not invalidate it for the easy regime.

\textbf{Translation to our setting.} Our noise is i.i.d.\ across every axis
--- time, neuron, simulation seed, stimulus frame. This means \emph{every}
axis is a valid permutation axis for an N2N construction, and we have more
flexibility than SPEND, not less. The challenge for us is not finding a
decorrelating axis (any axis works); it is choosing the axis whose
\emph{signal} is the most highly correlated, so the network has the
strongest signal to learn while the noise pair has independent realisations.

\subsection{What the signal-correlation structure of flyvis lets us exploit}

\begin{table}[h]
\centering
\small
\begin{tabular}{lll}
\toprule
\textbf{Domain} & \textbf{Signal correlation} & \textbf{Noise correlation} \\
\midrule
Time $t$ & strong (smooth $\sim$ms dynamics, $\Delta t = 20$ms) & uncorrelated \\
Cell-type $c_i$ (65 classes) & strong (typed equivariance) & uncorrelated \\
Spatial column (217 retinotopic columns) & strong (same RF, shifted stimulus) & uncorrelated \\
Stimulus replay (re-run with same $I(t)$, different seed) & identical (deterministic ODE) & uncorrelated \\
\bottomrule
\end{tabular}
\end{table}

For SPEND, axis selection was forced (only the spectral axis decorrelated
noise). For us, axis selection is \emph{strategic}: we pick the axis along
which the signal is strongest, because the N2N convergence rate is set by
how cleanly the network can predict signal from a noisy pair. The
stimulus-replay axis is the cleanest (deterministic signal, fully
independent noise), the time axis is the most data-efficient
(no extra simulation needed), and the cell-type/columnar axis is the most
mechanistically informative (it ties recovery to the connectome's structural
priors).

\textbf{One-line argument the user can use:} \emph{N2N requires only that
the two views share the same signal and have independent noise; i.i.d.\
Gaussian noise satisfies this in every axis. SPEND's permutation trick
generates such pairs from a single recording. We can use the same trick
along time, cell-type, or stimulus-replay axes --- the i.i.d.\ assumption
is what makes our problem easier, not harder, than SPEND's.}

\section{Three add-ons --- implementation plan}

Each add-on is a self-contained training-time intervention that can be
toggled by a single YAML coefficient. None modifies the GNN architecture.
All three address attack (A) --- reformulating the loss so it operates on a
denoised observable rather than on a noisy finite difference.

\subsection{Add-on \#1 --- N2N-on-time (frame-permutation training)}

\textbf{Mechanism.} Within each training window of $T_w$ consecutive frames
(say $T_w = 16$, $\Delta t = 20$\,ms, total 320\,ms --- well within the
correlation length of the dynamics), construct two views:
\[
v^{\text{even}}_{i,k} = v^{\text{obs}}_i(t_0 + 2k \Delta t), \quad
v^{\text{odd}}_{i,k} = v^{\text{obs}}_i(t_0 + (2k{+}1)\Delta t),\quad k=0,\ldots,T_w/2 - 1.
\]
The two views share the same underlying clean trajectory (subject to the
half-frame shift, which is small at this $\Delta t$) but have independent
noise realisations because $\xi_i(t) \perp \xi_i(t')$ for $t \ne t'$.

\textbf{Loss.} Train an auxiliary 1D-conv \emph{voltage smoother}
$\mathcal D: v^{\text{even}} \mapsto \widehat v^{\text{clean}}$ with the
N2N loss
$\mathcal L_{\text{N2N}} = \| \mathcal D(v^{\text{even}}) - v^{\text{odd, interp}}\|_2^2$
where $v^{\text{odd, interp}}$ is the half-frame interpolant of the odd
trace evaluated at even-frame timestamps. The cleaned trace
$\widehat v^{\text{clean}}$ then replaces $v^{\text{obs}}$ everywhere it
appears in the GNN forward pass and in the supervised target
$\widehat y = (\widehat v[t{+}1] - \widehat v[t])/\Delta t$.

\textbf{Why it helps.} The derivative noise std drops from
$\gamma\sqrt 2/\Delta t \approx 7.07$ to
$\sim \gamma/\sqrt{T_w \cdot \Delta t}$ --- order-of-magnitude SNR gain. The
neurips paper's pointwise MSE is replaced by an MSE on a target whose
variance has been reduced $\sim 8\times$, exactly SPEND's reported gain.

\textbf{Wire-up.} New regularizer component
\texttt{coeff\_n2n\_time\_loss} in
\texttt{src/connectome\_gnn/models/regularizer.py}; new module
\texttt{src/connectome\_gnn/models/voltage\_smoother.py} (10--20 LoC, 1D
U-Net with 3 levels, channels 16/32/64, total $\sim 30\text{k}$ params); call site
in \texttt{graph\_trainer.py} (3 lines). Smoother is co-trained with the
GNN; the N2N loss flows through it but not through $\widehat W$.

\textbf{Falsifiable PASS criterion (Phase~S).} On a held-out 16-frame window
of \texttt{flyvis\_noise\_free} corrupted with $\gamma=0.10$ Gaussian noise,
the smoother's output must reduce per-pixel std by $\geq 4\times$ versus
the input, while preserving signal mean to within 5\%.

\textbf{Risk.} The N2N target needs sub-frame interpolation; at
$\Delta t = 20$\,ms the dynamics changes by $\lesssim 5\%$ between adjacent
frames so linear interpolation is sufficient. If the smoother over-smooths,
$\widehat v$ loses high-frequency content needed by $g_\phi(v)^2$ (the
squaring is sensitive to fast voltage transients). Mitigation: regularize
the smoother's bandwidth via a small L2 on its second-derivative output.

\subsection{Add-on \#2 --- N2N-across-cells (typed-equivariance pair)}

\textbf{Mechanism.} For each cell type $c$ (65 total), pick two neurons
$i_1, i_2$ of that type that are exposed to approximately the same stimulus
(same retinotopic column with a small spatial offset, or successive frames
of a slow video stimulus). Their clean voltages should be
\emph{equivariant under stimulus shift}, i.e.\ $v_{i_1}(t) \approx v_{i_2}(t + \tau_{12})$
where $\tau_{12}$ is a known column-offset delay. After delay-correction,
the two traces share the same signal but have independent noise.

\textbf{Loss.}
$\mathcal L_{\text{equiv}} = \sum_{c,\,(i_1,i_2)\in c}
\|v^{\text{obs}}_{i_1}(t) - v^{\text{obs}}_{i_2}(t + \tau_{12})\|_2^2 - 2\sigma_\eta^2$
where the constant subtraction cancels the noise-variance bias so the
expected loss is exactly the residual cell-type heterogeneity.

\textbf{Why it helps.} This injects two pieces of information at once:
(i) typed equivariance breaks the \emph{continuous columnar degeneracy} of
\S\,1.1 directly, because the loss is non-zero on sum-zero columnar
perturbations; (ii) it provides an N2N-style noise reduction because the
loss is invariant under the mean-zero noise on $v_{i_1} - v_{i_2}$.

\textbf{Wire-up.} New regularizer component
\texttt{coeff\_typed\_equiv}; requires a precomputed lookup table
$(c,\, i_1,\, i_2,\, \tau_{12})$ from the connectome metadata
(already available via cell-type and column index). Cost is
$O(65)$ pairs per minibatch, negligible.

\textbf{Falsifiable PASS criterion.} On clean traces ($\gamma=0$), the loss
must be $< 10^{-3}$ at all timesteps for at least 60/65 cell types
(some types have only one column-instance and are excluded). Under
$\gamma=0.10$, the noise-cancelled estimator (subtracting $2\sigma_\eta^2$)
must remain unbiased to within 5\% of the clean value.

\textbf{Risk.} Columnar offsets are not perfectly known --- the connectome
gives us cell-type labels, but the column-index map has residual jitter from
the FlyWire registration. Mitigation: limit pairs to within-column same-type
pairs first (where $\tau_{12} = 0$), and add cross-column pairs in a
later iteration.

\subsection{Add-on \#3 --- Stimulus-replay residual denoising}

\textbf{Mechanism.} Generate two simulations with \emph{the same stimulus
$I(t)$} but \emph{different noise seeds}:
$v_a, v_b$ both $= v^{\text{clean}}(I) + \gamma\,\xi_{a,b}$. The
\emph{stimulus-conditional residual} $r_a = v_a - \mathbb E[v|I]$ is pure
noise; $r_b$ likewise. But $\mathbb E[v|I] = v^{\text{clean}}$ (deterministic
ODE), so $r_a + r_b = $ pure independent noise of zero mean. The N2N pair
is $(v_a, v_b)$ directly --- both views share the same clean signal
$v^{\text{clean}}$ and have independent noise.

\textbf{Loss.} A denoiser $\mathcal D: v_a \mapsto \widehat v^{\text{clean}}$
with $\mathcal L = \|\mathcal D(v_a) - v_b\|_2^2$. This is the closest
direct port of SPEND --- the ``permutation'' is replaced by a second
simulation seed.

\textbf{Why it helps.} Same N2N rationale as add-on \#1, but the two views
share the \emph{exact} clean signal (no half-frame shift) so there is no
interpolation loss. Cleanest of the three add-ons; cost is one extra
simulation pass per training epoch.

\textbf{Wire-up.} New simulation flag
\texttt{simulation.replay\_seed\_pair: true} in
\texttt{src/flyvis\_gnn/generators/graph\_data\_generator.py}; new dataset
field \texttt{v\_obs\_pair} (shape $T \times N \times 2$); new regularizer
component \texttt{coeff\_replay\_n2n}.

\textbf{Falsifiable PASS criterion.} Under
$\gamma=0.10$, the trained denoiser's output on a held-out
stimulus must achieve $\|\mathcal D(v_a) - v^{\text{clean}}\|_2 / \|v^{\text{clean}}\|_2
< 0.04$ (4\% relative error), versus 10\% for the raw $v_a$.

\textbf{Risk.} Doubles simulation cost. Mitigation: re-use the same DAVIS
clip with two noise realisations; this only doubles memory, not
compute, since the clean ODE rollout is shared.

\subsection{Recommended order and integration with the agentic loop}

\begin{enumerate}
  \item \textbf{Add-on \#3 first} (stimulus-replay) --- cleanest mathematics,
    smallest implementation surface, sets a ceiling on what N2N can buy us.
    If add-on \#3 fails to lift $W$ R$^2$, no other N2N variant will help.
  \item \textbf{Add-on \#1 second} (time-permutation) --- gets the same
    benefit without doubling simulation cost; the ratio of its result to
    add-on \#3's tells us how much we lose from the half-frame-shift
    approximation.
  \item \textbf{Add-on \#2 last} (typed equivariance) --- this one
    \emph{also} attacks the symmetry directly (anchor class~4 in the
    instruction-file taxonomy), so it is independent of \#1/\#3 and worth
    running in parallel.
\end{enumerate}

Each add-on lives in a single block of the agentic loop, conforming to the
existing Phase R/S/C/V protocol
(\texttt{instruction\_flyvis\_noise\_005\_010\_code\_change.md:373--559}). The
PASS criteria above plug directly into Phase~S; the verdict criteria
($\Delta W$~R$^2 \geq 0.005$, $\geq 75\%$ seeds better, no catastrophe)
are unchanged.

\section{Practical takeaway}

The bottleneck under $\gamma=0.10$ is not the GNN architecture --- it is the
combination of an \emph{absorbing} architectural symmetry
(W~$\leftrightarrow f_\theta$ scale) with a \emph{symmetry-blind} loss
(pointwise MSE on a noisy finite difference). Voltage-level reweighting
cannot move the optimizer off the slope-shrunk minimum. SPEND validates a
broader principle: \emph{reformulating the loss along an axis where noise
decorrelates dissolves the bottleneck}. The fact that SPEND was designed for
non-independent noise is irrelevant to whether it transfers --- the
underlying Noise2Noise theorem requires only independence between two
noise realisations, which i.i.d.\ Gaussian satisfies trivially. Three concrete
add-ons (stimulus-replay, time-permutation, typed-equivariance) are
implementable as single-block agentic interventions and address attack (A)
of the bottleneck head-on.

\end{document}
